\documentclass[12pt,a4paper,french]{book}
\usepackage[utf8]{inputenc}
\usepackage[T1]{fontenc}
%\usepackage[francais]{babel}
\usepackage{polyglossia}
\setmainlanguage{french}
\setmainlanguage{english}
\usepackage{graphicx}
\usepackage{caption}
\usepackage{subcaption}
\usepackage{cprotect}
\usepackage{geometry}
\usepackage{graphicx}
\usepackage{xcolor}
\usepackage{hyperref}
\usepackage{fancyhdr}
\usepackage{array}
\usepackage{booktabs}
\usepackage{listings}
\usepackage{xhfill}
\usepackage{amsmath}
\usepackage{amssymb}
\usepackage{textcomp}
\usepackage{float}
\usepackage{verbatim}
\usepackage{titlesec}
\usepackage{cprotect}
\usepackage{float}
\usepackage{bm}
\usepackage{forest}
% Géométrie de la page
\geometry{
    left=1.0cm,
    right=1.0cm,
    top=.5cm,
    bottom=1.5cm,
    headheight=1.0cm,
    headsep=1cm,
    footskip=0.75cm
}

% Couleurs
\definecolor{darkblue}{HTML}{1f4788}
\definecolor{lightblue}{HTML}{E8F0F8}
\definecolor{darkgreen}{HTML}{2d7d3a}
\definecolor{lightgreen}{HTML}{E8F5E9}
\definecolor{darkred}{HTML}{c1121f}
\definecolor{lightyellow}{HTML}{FFF8DC}
\definecolor{codegray}{gray}{0.95}
\definecolor{lightcyan}{HTML}{E0F7FA}
\definecolor{darkcyan}{HTML}{00838F}
% En-tête et pied de page
\pagestyle{fancy}
\fancyhf{}
\fancyhead[L]{\textcolor{darkblue}{\textbf{PyCyto\_Emergence}}}
\fancyhead[R]{\textcolor{darkblue}{Guide d'utilisation}}
\fancyfoot[C]{\thepage}
\fancyfoot[R]{\small v0.6.0}
\renewcommand{\headrulewidth}{1pt}
\renewcommand{\footrulewidth}{0.5pt}

% Styles des sections
\titleformat{\chapter}[display]
{\normalfont\huge\bfseries\color{darkblue}}
{\chaptertitlename\ \thechapter}
{20pt}
{\huge}

\titleformat{\section}
{\normalfont\Large\bfseries\color{darkblue}}
{\thesection}
{1em}
{}

\titleformat{\subsection}
{\normalfont\large\bfseries\color{darkgreen}}
{\thesubsection}
{1em}
{}

% Environnement pour les boîtes
\usepackage[most]{tcolorbox}

\newenvironment{infobox}
{\begin{tcolorbox}[colback=lightblue, colframe=darkblue, title=\textbf{ℹ Information}, boxrule=1.5pt]}
{\end{tcolorbox}}

\newenvironment{successbox}
{\begin{tcolorbox}[colback=lightgreen, colframe=darkgreen, title=\textbf{✓ Succès}, boxrule=1.5pt]}
{\end{tcolorbox}}

\newenvironment{warningbox}
{\begin{tcolorbox}[colback=lightyellow, colframe=darkred, title=\textbf{⚠ Attention}, boxrule=1.5pt]}
{\end{tcolorbox}}


\newenvironment{todobox}
{\begin{tcolorbox}[colback=lightcyan, colframe=darkcyan, 
		title=\textbf{[>] To Do}, boxrule=1.5pt]}
{\end{tcolorbox}}
% Code listings
\lstset{
    basicstyle=\ttfamily\small,
    breaklines=true,
    backgroundcolor=\color{codegray},
    frame=single,
    frameround=tttt,
    captionpos=b,
    tabsize=4,
    language=Python
}




\usepackage{tikz}
\usetikzlibrary{arrows.meta}
\tikzset{%
	>={Latex[width=2mm,length=2mm]},
	% Specifications for style of nodes:
	base/.style = {rectangle, rounded corners, draw=black,
					minimum width=4cm, minimum height=1cm,
					text centered, font=\sffamily},
	algorithmStarts/.style = {base, fill=blue!30},
	startstop/.style = {base, fill=red!30},
	function/.style = {base, fill=green!30},
	process/.style = {base, minimum width=2.5cm, fill=orange!15,font=\ttfamily},
}


% Titre personnalisé
\title{\Huge\textbf{\textcolor{darkblue}{ATOMOD}}\\
       \Large\textcolor{darkgreen}{User guide}\\
       }
\author{H. Bulou}
\date{Version 0.6.0 \\ \today}

% ============================================================================
% DOCUMENT
% ============================================================================

\begin{document}

% Page de titre
\maketitle
% Page blanche
\newpage
\thispagestyle{empty}
\mbox{}

% Table des matières
%\newpage
%\tableofcontents
%\newpage
% ============================================================================
% INSTALL
% ============================================================================
\section*{Installation (linux)}

%AVANT DE COMMENCER : CONFIGURATION REQUISE POUR ATOMOD  ⚠
%Afin d'éviter tout échec en cours de route, assurez-vous de disposer de :
% ESPACE DISQUE :
%   • Au moins 7 GB d'espace libre sur votre partition.
% LOGICIELS SYSTÈME (Obligatoires) :
%   • git          (pour cloner le dépôt ATOMOD)
%   • python3-venv (pour isoler les paquets Python)
%   • python3-pip  (pour installer les dépendances)
%   • wget         (pour récupérer les fichiers du tutoriel)
% Sur Ubuntu / Xubuntu / Debian, vous pouvez tout installer avec :
% sudo apt update && sudo apt install git python3-pip python3-venv wget


Before you begin, check the ATOMOD systems requirements, to avoid any issues during the process.
Ensure you have the following:
\begin{itemize}
	\item DISK SPACE:
	\begin{itemize}
		\item 	At least 7 GB of free space on your partition. 
	\end{itemize}
	\item SYSTEM SOFTWARE (Required):
	\begin{itemize}
		\item \verb+git+          (to clone the ATOMOD repository)
		\item \verb+python3-venv+ (to isolate Python packages)
		\item \verb+python3-pip+  (to install dependencies)
		\item  \verb+wget+         (to download installer \verb+install.py+)
	\end{itemize}
	On Ubuntu / Xubuntu / Debian, you can install everything using:\\
	\verb+sudo apt update && sudo apt install git python3-pip python3-venv wget+
	
\end{itemize}

Installation process on linux:
\begin{itemize}
	\item 1 - Open an X terminal
	\item 2 - Navigate to the directory of your choice (e.g., \verb+cd ~/Desktop+ or \verb+cd ~/src+)(Fig.\ref{fig:install1}(a)).
	
	\verb+cd Desktop+
	
	
	
	\item 3 - Download the installation script \verb+install.py+ (Fig.\ref{fig:install1}(b)).
	
	\verb+wget https://github.com/hbulou/ATOMOD/releases/download/v0.1/install.py+
	
	\item 4 - Run the installation (Fig.\ref{fig:install1}(c)). NOTE: You must have at least 7 GB of available disk space. 
	
	\verb+python3 install.py+
	
	\item 5 - Once the installation is complete, go to the ATOMOD directory (Fig.\ref{fig:install1}(d)).
	
	\verb+cd ATOMOD+
	
	\item 6 - Run the \verb+run-tuto-ATOMOD.sh+ script to launch the tutorials (Fig.\ref{fig:install1}(e)).
	
	\verb+./run-tuto-ATOMOD.sh+
	
\end{itemize}

\begin{figure}[H]
	\centering
	\includegraphics[width=1.0\linewidth]{install1}
	\caption{Screenshots of the various stages of installing ATOMOD on Linux.}
	\label{fig:install1}
\end{figure}


% ============================================================================
% INDEX
% ============================================================================

%\printindex
\bibliographystyle{plain} 
%\bibliography{ATOMOD}
%\begin{thebibliography}{3}
\end{document}
